% !PRINTED 395
\weberpage{395}
9. \gesperrt{If, then, we denote by $\nu$ the number of the different odd primes dividing $D$, then by \S\,98, 12. we have}
\[
\begin{aligned}
(Z, R) &= 2^{\nu}, & D &\equiv 1 \ (\mathrm{mod}\ 4), \quad D \equiv 4 \ (\mathrm{mod}\ 16),\\
(Z, R) &= 2^{\nu+1}, & D &\equiv 8, 12, 16, 24, 28 \ (\mathrm{mod}\ 32),\\
(Z, R) &= 2^{\nu+2}, & D &\equiv 0 \ (\mathrm{mod}\ 32),
\end{aligned}
\]
\gesperrt{or in general}
\[
(Z, R) = 2^{\lambda},
\]
\gesperrt{where $\lambda$ has the same meaning as in \S\,104, 3.}

This number agrees, as one sees, with the number of the fundamental divisors of $D$, and from 8. it follows that an absolute norm residue $r$ retains this property when $r$ is increased by a multiple of $D$.

From this the theorem:

10. \gesperrt{The absolute norm residues are contained in a finite number of arithmetic progressions which advance by multiples of $D$.}

The odd primes dividing $D$, and in cases a) and b) of No. 8 the number 8 or 4, are called the \gesperrt{characteristic primes and prime powers of the order}.

\begin{center}
\S\,109. \gesperrt{The Genera of Ideals.}
\end{center}

According to the foregoing we can decompose the totality of the rational numbers $z$, excluding those which are not relatively prime to $D$ in the numerator or in the denominator, into $\mu$ classes (subsidiary groups), where
\weq{1}{(Z, R) = \mu}
is the value determined above. Accordingly we have:
\weq{2}{Z = R + R_1 + R_2 \ldots + R_{\mu - 1}.}
If a number $z$ belongs to one of these subsidiary groups, then all numbers congruent to $z$ with respect to the modulus $D$ belong to the same subsidiary group.

11. \gesperrt{The principal class $R$ of the norm residues is a group with respect to multiplication and division.}

If $r$ belongs to the principal class $R$, then two numbers $z$ and $rz$ belong to \gesperrt{one and the same class}.

The totality of the integral and fractional ideals $\mathfrak{a}$ of the field $\Omega$ that are coprime to $D$ also form a group $\overline{O}$.
